\section{Resumen ejecutivo}

\pendiente{Redactar entre 300 y 400 palabras. Contar al final con un script; el límite es parte del
requisito.}

% ---------------------------------------------------------------------
% Guía de redacción (borrar los comentarios al terminar):
%  1. Qué problema real resuelve el sistema y para qué organización
%     concreta (nombre o seudónimo autorizado en el acta), no "las empresas".
%  2. Quién lo usará: clientes, mesa de ayuda, técnicos de campo.
%  3. Cómo se resuelve hoy el problema y qué consecuencia medible tiene
%     (tiempo, costo, errores) según el acta de la entrevista.
%  4. Por qué la solución debe ser distribuida: dos razones ligadas
%     directamente a RNF-01 (disponibilidad) y RNF-02 (escalabilidad),
%     sin definiciones de manual.
%  5. Alcance del periodo: qué se construye en E1, E2, E3 y E4, y qué
%     queda explícitamente fuera.
% ---------------------------------------------------------------------

\section{Contexto y problema}

\subsection{Organización y usuarios}

\pendiente{Describir la organización real (o el grupo de usuarios) con los datos del acta:
nombre o seudónimo autorizado, actividad, tamaño, ubicación, horario de atención y volumen de
casos por semana. Si se usa seudónimo, indicarlo y explicar por qué.}

\subsection{Cómo se resuelve hoy}

\pendiente{Describir el proceso actual con sus herramientas reales (hoja de cálculo, WhatsApp,
cuaderno, sistema anterior), quién hace qué y en qué orden. Un párrafo por paso del proceso.}

\subsection{Consecuencias medibles de no resolverlo}

Cada cifra de esta tabla debe poder rastrearse hasta una frase del acta (anexo~\ref{anexo:acta}).

\begin{table}[H]
\centering
\caption{Consecuencias actuales del problema}
\label{tab:consecuencias}
\small
\begin{tabularx}{\linewidth}{@{}lYY@{}}
\toprule
Indicador & Valor actual & Fuente \\
\midrule
Tiempo medio de registro de una incidencia & \pendiente{cifra} & Acta, pregunta \pendiente{n} \\
Tiempo medio de asignación al técnico & \pendiente{cifra} & Acta, pregunta \pendiente{n} \\
Casos perdidos o duplicados por semana & \pendiente{cifra} & Acta, pregunta \pendiente{n} \\
Incidencias que vencen el SLA sin aviso & \pendiente{cifra} & Acta, pregunta \pendiente{n} \\
Costo estimado por hora de servicio detenido & \pendiente{cifra} & Acta, pregunta \pendiente{n} \\
\bottomrule
\end{tabularx}
\end{table}

\subsection{Interesados y expectativas}

\begin{table}[H]
\centering
\caption{Interesados del sistema y lo que esperan}
\label{tab:stakeholders}
\small
\begin{tabularx}{\linewidth}{@{}lYY@{}}
\toprule
Interesado & Qué espera del sistema & Cómo se comprueba \\
\midrule
Cliente final & Saber en qué va su caso sin llamar & Consulta de estado con número de ticket \\
Agente de mesa de ayuda & Registrar y clasificar rápido, sin duplicar & Formulario de alta en menos de 60 s \\
Técnico de campo & Ver su cola y actualizar el estado desde el celular & Interfaz adaptable y actualización sin conexión estable \\
Responsable de operaciones & Cumplir el SLA y ver la carga por técnico & Reporte de cumplimiento y tablero \\
Gerencia & Reducir el costo de atención & Indicador mensual de costo por caso \\
\bottomrule
\end{tabularx}
\end{table}

\subsection{Alcance del periodo}

\begin{table}[H]
\centering
\caption{Alcance por entrega}
\small
\begin{tabularx}{\linewidth}{@{}lY@{}}
\toprule
Entrega & Alcance comprometido \\
\midrule
E1 & Definición, requisitos, arquitectura C4, ADR, CAP/PACELC y esqueleto de dos servicios \\
E2 & Microservicios completos, API documentada y despliegue en contenedores \\
E3 & Base de datos distribuida y procesamiento de datos del SLA \\
E4 & Sistema final en capas, pruebas, calidad y defensa \\
\midrule
Fuera del alcance & \pendiente{facturación, aplicación móvil nativa, integración con el sistema del proveedor de red} \\
\bottomrule
\end{tabularx}
\end{table}

\section{Objetivos}

\subsection{Objetivo general}

Construir el SGIRI como un sistema distribuido que permita a un proveedor de internet registrar
cualquier incidencia técnica que reporten sus clientes y asignarla a un técnico especializado en
ese tipo de problema, de modo que una falla en una parte del sistema no detenga el registro ni la
atención de las incidencias.

\subsection{Objetivos específicos e indicadores}

\begin{table}[H]
\centering
\caption{Objetivos específicos y su indicador de cumplimiento}
\label{tab:objetivos}
\small
\begin{tabularx}{\linewidth}{@{}clY@{}}
\toprule
\# & Objetivo específico & Indicador con que se comprueba \\
\midrule
OE1 & Diseñar la arquitectura y comunicar dos servicios (E1)
& Los servicios de incidencias y técnicos responden por REST y devuelven 503 cuando el otro está
apagado; se comprueba con capturas el 29/11/2026 \\
OE2 & Agregar el API Gateway, las notificaciones con RabbitMQ y el despliegue con Docker Compose (E2)
& El sistema arranca con un solo \texttt{docker compose up} y las pruebas de humo pasan en la
semana 9 \\
OE3 & Probar la tolerancia a fallos y la carga (E3)
& Con 200 usuarios simultáneos, el 95\,\% de las consultas responde en menos de 300~ms (RNF-03 y
RNF-04), medido en la semana 13 \\
OE4 & Entregar el sistema completo con todos los requisitos funcionales (E4)
& Los 11 requisitos funcionales pasan sus criterios de aceptación y la disponibilidad medida
llega al 99,5\,\% (RNF-01) en la semana 17 \\
\bottomrule
\end{tabularx}
\end{table}
